% macros.tex — preamble and macros for the OpenBench SRS (D17).
%
% Requirement identifiers
%   \fr{Inventory.Type.Create}{The system shall ...}   functional requirement, dot-notation ID, label fr:<ID>
%   \fr[Stretch]{Metered.Tier.Define}{...}             optional tag printed after the ID
%   \FR{Inventory.Type.Create}                          hyperlinked reference to a functional requirement
%   \br{key}{text}  /  \BR{key}                         business rule BR-n and its reference
%   \fe{key}{text}  /  \FE{key}                         major feature FE-n and its reference
%   Other numbered classes follow the same pattern: \bo \scc \li \ri \co \ud \asm \dep \ui \si \ci \pe \se \qa
%   and their references \boref \sccref \liref ... (see \newidclass below).
% Use cases (full use case table)
%   \begin{usecase}{key}{Name} \ucfield{Actors}{...} ... \end{usecase}   and \UC{key}

\usepackage[letterpaper,margin=1in]{geometry}
\usepackage[T1]{fontenc}
\usepackage{lmodern}
\usepackage{microtype}
\usepackage{parskip}
\usepackage{enumitem}
\usepackage{booktabs}
\usepackage{longtable}[=v4.13]   % v4.24 (TeX Live 2025) raises "Infinite glue shrinkage" on every multi-page table; v4.13 does not
\usepackage{array}
\usepackage{tabularx}
\usepackage{multirow}
\usepackage{makecell}
\usepackage{multicol}
\usepackage{pdflscape}   % landscape pages for the wide figures (state diagram, data model)
\usepackage{xcolor}
\usepackage{graphicx}
\usepackage{caption}
\usepackage{fancyhdr}
\usepackage{lastpage}
\usepackage{needspace}
\usepackage{etoolbox}
\usepackage[colorlinks=true,linkcolor=blue!50!black,urlcolor=blue!50!black,citecolor=blue!50!black]{hyperref}

% ---- Document metadata -------------------------------------------------------
\newcommand{\SRSversion}{0.3}
\newcommand{\SRSdate}{2026-10-05}      % date of this version
\newcommand{\SRScreated}{2026-09-19}   % date of version 0.1 (use case creation date)
\newcommand{\SRSproduct}{OpenBench}

% ---- Layout ------------------------------------------------------------------
\setcounter{secnumdepth}{3}
\setcounter{tocdepth}{2}
\setlist{itemsep=1pt,topsep=2pt,parsep=0pt}
\setlist[enumerate]{leftmargin=*}
\setlist[itemize]{leftmargin=*}
\captionsetup{font=small,labelfont=bf}
\graphicspath{{build/figures/}}
\newcolumntype{L}[1]{>{\raggedright\arraybackslash}p{#1}}
\newcolumntype{Y}{>{\raggedright\arraybackslash}X}

\pagestyle{fancy}
\fancyhf{}
\lhead{\small \SRSproduct{} Software Requirements Specification}
\rhead{\small Version \SRSversion}
\cfoot{\small Page \thepage\ of \pageref*{LastPage}}
\renewcommand{\headrulewidth}{0.4pt}
\fancypagestyle{plain}{\fancyhf{}\cfoot{\small Page \thepage\ of \pageref*{LastPage}}\renewcommand{\headrulewidth}{0pt}}

% Vertical gap after each numbered item (\fr, \br, \glossterm ...); the short edition sets it smaller.
\newlength{\reqgap}\setlength{\reqgap}{3pt}

% ---- Generic numbered identifier classes --------------------------------------
% \newidclass{co}{CO} defines \co{key}{text} (prints "CO-n  text", labels co:key)
% and \coref{key} (hyperlinked "CO-n").
\newcommand{\newidclass}[2]{%
  \newcounter{#1}%
  \expandafter\newcommand\csname #1\endcsname[2]{%
    \par\noindent\refstepcounter{#1}\label{#1:##1}%
    \hangindent=5em\hangafter=1%
    \makebox[5em][l]{\textbf{#2-\arabic{#1}}}##2\par\vspace{\reqgap}}%
  \expandafter\newcommand\csname #1ref\endcsname[1]{\hyperref[#1:##1]{#2-\ref*{#1:##1}}}%
}
\newidclass{bo}{BO}    % business objective
\newidclass{scc}{SC}   % success criterion
\newidclass{fe}{FE}    % major feature
\newidclass{li}{LI}    % limitation / exclusion
\newidclass{ri}{RI}    % business risk
\newidclass{co}{CO}    % design and implementation constraint
\newidclass{ud}{UD}    % user documentation
\newidclass{asm}{AS}   % assumption
\newidclass{dep}{DE}   % dependency
\newidclass{ui}{UI}    % user interface
\newidclass{si}{SI}    % software interface
\newidclass{ci}{CI}    % communications interface
\newidclass{pe}{PE}    % performance requirement
\newidclass{se}{SE}    % security requirement
\newidclass{qa}{QA}    % quality attribute
\newidclass{br}{BR}    % business rule
\newcommand{\FE}{\feref}
\newcommand{\BR}{\brref}

% ---- Functional requirements (dot-notation IDs) --------------------------------
\newcommand{\fr}[3][]{%
  \par\noindent\phantomsection\label{fr:#2}%
  \hangindent=2em\hangafter=1%
  \textbf{#2}%
  \if\relax\detokenize{#1}\relax\else\ \textcolor{black!60}{[#1]}\fi
  \quad #3\par\vspace{\reqgap}}
\newcommand{\FR}[1]{\hyperref[fr:#1]{#1}}

% ---- Use cases (full use case table) -------------------------------------------
\newcounter{uc}
\newenvironment{usecase}[2]{%
  \refstepcounter{uc}\label{uc:#1}%
  \needspace{10\baselineskip}%
  \subsection*{UC-\theuc: #2}%
  \addcontentsline{toc}{subsection}{UC-\theuc: #2}%
  \begin{longtable}{@{}>{\bfseries\raggedright\arraybackslash}p{0.2\textwidth}>{\raggedright\arraybackslash}p{0.76\textwidth}@{}}
  \toprule
  Use Case ID & UC-\theuc \\
  Use Case Name & #2 \\
  \endfirsthead
  \multicolumn{2}{@{}l}{\small\itshape UC-\theuc\ #2 (continued)}\\ \toprule
  \endhead
  \midrule \multicolumn{2}{@{}r}{\small\itshape continues on next page}\\
  \endfoot
  \bottomrule
  \endlastfoot
}{\end{longtable}}
\newcommand{\ucfield}[2]{#1 & #2 \\}
\newcommand{\ucflow}[2]{\textbf{#1\ #2}\par}
\newcommand{\UC}[1]{\hyperref[uc:#1]{UC-\ref*{uc:#1}}}

% ---- Glossary entries (emitted by tools/gen_glossary.py) ----------------------
\newcommand{\glossterm}[3]{%
  \par\noindent\hangindent=1.5em\hangafter=1%
  \textbf{#1}\quad #2%
  \if\relax\detokenize{#3}\relax\else\ {\small\textit{Avoid:} #3.}\fi\par\vspace{\reqgap}}

% Condensed use case (short SRS): a two-column table whose number comes from the label \ucrow set in the summary table.
% The title is the table's first row; rows are \ucfield{Label}{text} as in the full template.
% \ucrow{key}{Name}{Actor}{Trigger and outcome}: one row of the use case summary table (short edition); defines uc:key.
\newcommand{\ucrow}[4]{\refstepcounter{uc}\label{uc:#1}UC-\theuc & #2 & #3 & #4 \\}
\newcommand{\glossgroup}[1]{\par\medskip\noindent\textbf{#1}\par\smallskip}

% ---- Small helpers -------------------------------------------------------------
\newcommand{\priority}[1]{\textbf{Priority:} #1}
\newcommand{\brgroup}[1]{\par\medskip\noindent\textbf{#1}\par\nobreak\smallskip}   % heading for a group of business rules
\newcommand{\term}[1]{\emph{#1}}                 % first use of a glossary term in a section
\newcommand{\stretchtag}{\textcolor{black!60}{[Stretch]}}
\newcommand{\figref}[1]{Figure~\ref{#1}}
